%% Chapter 10: Exact Exponential-Time Algorithms I: Branching
%% Status: skeleton.  Only the learning goals are fixed; they come from the
%% exam pool.  The sectioning is deliberately not invented in advance --
%% it will follow the source material once that has been read.

\chapter{Exact Exponential-Time Algorithms I: Branching}
\label{ch:exact-i}

If a problem is NP-hard and we still need the exact optimum, the question is no
longer polynomial versus exponential but $2^n$ versus $1.5^n$ --- a difference of
a factor of a million already at $n = 40$. Branching, together with reduction
rules that shrink the instance, is the basic technique.

\begin{goals}
  \poolgoal{use $\Ostar$ notation and explain why polynomial factors may be suppressed}
  \poolgoal{design branching algorithms and analyse them with branching vectors and recurrences}
  \poolgoal{for \lang{Exact 3-Satisfiability}, give a polynomial-time reduction rule transforming an instance into an equivalent smaller one in which all clauses have size two or three}
  \poolgoal{for the same problem, give a reduction rule after which all clauses have size exactly three}
  \poolgoal{give an $\Ostar(c^n)$ algorithm with $c < 2$ for it, and analyse its running time}
  \poolgoal{give an $\Ostar(c^n)$ algorithm with $c < 2$ for \lang{Bisection Width}}
\end{goals}

\todo{This chapter is not written yet. Write it against
\texttt{material/slides/Exponential Time Algorithms 2026.pptx}.}
